\section{组织、产品与中美差异}

访谈后半段从技术转向组织和产品。姚顺宇对 neo lab 潮流并不乐观，认为大多数都会死，除非有真正好的人和清晰方向。谈到中国与美国的产品差异时，他认为美国过去十年更擅长 ToB、enterprise 和效率软件，因为这些市场直接且利润高；中国更擅长复杂 C 端产品，把广告、直播、电商、内容和分发转成隐性的利润飞轮。

\begin{figure}[H]
\centering
\includegraphics[width=0.88\textwidth]{frame-org-products.jpg}
\caption{组织与产品章节：讨论 neo lab、C 端产品、美国 enterprise 市场和中国复杂产品能力。\protect\footnotemark}
\end{figure}
\footnotetext{视频画面时间区间：03:12:57--03:13:28。该帧来自组织与产品讨论中段。}

\subsection{neo lab 的机会与死亡率}

很多从模型大厂出来的团队会成立新实验室，但姚顺宇认为多数 neo lab 会死。原因是语言模型主线已经不是蓝海，基础设施、算力、人才、数据、评估、产品入口都越来越集中。除非团队在新方向上有真正差异化，否则很难与大模型公司竞争。

\begin{knowledgebox}{为什么“语言模型末班车已发车”}
这句话不是说 AI 没机会，而是说最核心的语言模型主战场已经高度资本化、组织化和基础设施化。年轻研究者更适合寻找尚未被做好的方向，例如多模态生成、机器人、科学 AI、量子调控、长程 agent 等。
\end{knowledgebox}

\subsection{TPU 与 GPU 的组织含义}

访谈里有一段硬件讨论：GPU 的优势在生态和小规模通用性，典型节点内 NVLink 互连很强；TPU 则更像为大规模集群设计，采用 3D Torus 等拓扑。如果 compiler 和 sharding 策略足够好，可以利用更大规模的结构。这里的 sharding 指把参数、激活、优化器状态或数据切分到多张加速卡上，使它们不必在每张卡完整复制；它既是内存策略，也是通信策略。这里的关键不是哪个硬件绝对更好，而是硬件形态会影响软件栈、组织能力和训练策略。

\noindent\begin{tabularx}{\textwidth}{p{0.18\textwidth}p{0.35\textwidth}X}
\toprule
硬件路线 & 优势 & 代价 \\
\midrule
GPU & 生态成熟，通用性强，小规模易用；节点内高速互联 & 大规模成本高，供应链和集群调度压力大。 \\
TPU & 为大规模训练优化，拓扑可扩展，Google 内部软件栈深度配合 & 通用生态弱，更依赖 compiler、sharding 和内部基础设施。 \\
\bottomrule
\end{tabularx}

\begin{importantbox}{硬件不是孤立变量}
TPU/GPU 选择背后是组织能力选择：谁能写 compiler，谁能设计 sharding，谁能调度大集群，谁能把模型、数据和硬件一起优化。硬件优势必须通过软件栈和组织执行释放。
\end{importantbox}

\subsection{本章小结}

组织与产品讨论把技术竞争放回现实：模型公司不是论文团队，而是基础设施、产品、组织、市场和人才系统。neo lab 的生存、C 端产品的复杂度、硬件路线的选择，都取决于这种系统能力。

